% -*- coding: utf-8 -*-
\begin{thebibliography}{99}
\addcontentsline{toc}{chapter}{参考文献}

\bibitem{xie2021deep}
XIE H, QIN Z, LI G Y, et al. Deep learning enabled semantic communication systems[J]. IEEE Transactions on Signal Processing, 2021, 69: 2663--2675. DOI: 10.1109/TSP.2021.3071210.

\bibitem{dai2022nonlinear}
DAI J, WANG S, TAN K, et al. Nonlinear transform source-channel coding for semantic communications[J]. IEEE Journal on Selected Areas in Communications, 2022, 40(8): 2300--2316. DOI: 10.1109/JSAC.2022.3180802.

\bibitem{huang2025d2}
HUANG J, YUAN K, HUANG C, et al. D²-JSCC: Digital deep joint source-channel coding for semantic communications[J]. IEEE Journal on Selected Areas in Communications, 2025, 43(4): 1246--1261. DOI: 10.1109/JSAC.2025.3531546.

\bibitem{lyu2024semantic}
LYU Z, ZHU G, XU J. Semantic communications for image recovery and classification via deep joint source and channel coding[J]. IEEE Transactions on Wireless Communications, 2024, 23(8): 8388--8404. DOI: 10.1109/TWC.2023.3349330.

\bibitem{wang2023wireless}
WANG S, DAI J, LIANG Z, et al. Wireless deep video semantic transmission[J]. IEEE Journal on Selected Areas in Communications, 2023, 41(1): 214--229. DOI: 10.1109/JSAC.2022.3221977.

\bibitem{nguyen2025survey}
NGUYEN L X, RAHA A D, AUNG P S, et al. A contemporary survey on semantic communications: Theory of mind, generative AI, and deep joint source-channel coding[J]. IEEE Communications Surveys \& Tutorials, 2026, 28: 2377--2417. DOI: 10.1109/COMST.2025.3616973.

\end{thebibliography}
